\reportsection{Contributions}
\label{sec:contributions}

Merkle commitments are essential to modern proof systems, but developers often have to balance
performant systems with evolving security requirements and configurability, leaving underconfigured
schemes vulnerable. We aim to lift that burden into infrastructure.

\begin{itemize}[itemsep=0.35\baselineskip, topsep=0.35\baselineskip]
  \item \textbf{Main contribution}: we provide a rule for evaluating a Merkle commitment
  scheme's binding and hiding security against a target security parameter $\lambda$,
  field $\mathbb{F}$, hash function (or random oracle) $\mathcal{H}$, and hasher
  implementation (\cref{sec:background}).
  \item \textbf{Cost measurement}: we benchmark prover time, verifier time, and proof
  size across parameter combinations spanning three fields and two hash families,
  comparing desired security thresholds with their cost.
  \item \textbf{Machine checking our results}: we give the binding derivation formalised
  in Lean~4~\cite{ZLeanVC26}.
  \item \textbf{Batch opening consolidation}: we consolidate Arkworks' batch opening onto
  CoSet pruning, carrying forward the first semester's proof-size win over its
  prefix-pruning predecessor: roughly $360$\,kB down to roughly $33$\,kB at $4{,}096$
  opened leaves.
\end{itemize}

In two semesters of work we investigated the structure of Merkle trees and Merkle
commitments in depth, and what does and does not work in our favour.  Chronologically,
we started by looking at feature requests against the Merkle tree
implementation in Arkworks' \textit{crypto\_primitives}, which already sketched out many
missing features, evidenced by an RFC (\cref{app:rfc}).  We scoped
the problem to optimal batched index-proof pruning.  As a reference we used Chiesa and Yogev's
\emph{Building Cryptographic Proofs from Hash Functions}~\cite{CY26}, for its close
agreement with our definitions and the rigour of its proofs.  For batch proofs, the book
outlines a pruning algorithm we call \emph{CoSet pruning} (\cref{app:coset}).  Arkworks'
\textit{crypto\_primitives} used a distinct algorithm for the same objective, which we
call \emph{prefix pruning} (\cref{app:prefix}).  That comparison found no statistically
significant trade-off on prover or verifier time and significantly saved on proof size.

In the second semester we shifted to the proposed Arkworks design shift, aimed at
unifying tools that had previously been specific to individual constraint systems into
one IOR/polynomial IOP compiler (\cref{app:compilation}).  Our contribution there was scoped
to evaluating Vector Commitments and their Merkle-commitment subset for the hiding and
binding security they bring to the library.  A VC scheme must satisfy $\lambda$-bit
binding (\cref{sec:background}) and, optionally, $\lambda$-bit hiding
(\cref{sec:background}).  For a Merkle scheme, both properties are enforced through the
underlying hasher, and an undersized hasher still compiles and passes every functional
test, which is what makes the problem easy to miss.  Our research surfaced concerns that
the design shift had overlooked, and this report makes them explicit for developers and
users constructing a hasher for a target hiding and binding security level.
